%% ============================================================
%%  Lecture 2 -- Monday--Thursday, 13--16 July 2026
%%  Subfile of main.tex: compiles standalone OR as part of the book.
%%  Built from sources/2026-07-13/ (slide-transcript.md, outline.md).
%%  The deck ("2.1 Lecture", 18 slides) is agenda-style: most slides
%%  name what is developed at the board, and these notes work each
%%  named item in full (provenance tracked in the outline).
%% ============================================================
\documentclass[main.tex]{subfiles}

\begin{document}

\beginlecture{2}{Monday--Thursday, 13--16 July 2026}%
  {Estimation: Unbiasedness, Maximum Likelihood, and Confidence}

\epigraph{\ldots the object of statistical methods is the reduction of data.}%
{R. A. Fisher, ``On the Mathematical Foundations of Theoretical Statistics'' (1922)}

\begin{lecturecontents}{Contents of this lecture}
  \item \hyperlink{2.1}{2.1\quad Estimation as a Decision Problem}
  \item \hyperlink{2.2}{2.2\quad The Efficiency Bound, Revisited}
  \item \hyperlink{2.3}{2.3\quad Two Information Computations}
  \item \hyperlink{2.4}{2.4\quad Maximum Likelihood}
  \item \hyperlink{2.5}{2.5\quad Four Samples, Four Maximum-Likelihood Estimators}
  \item \hyperlink{2.6}{2.6\quad Several Parameters at Once}
  \item \hyperlink{2.7}{2.7\quad Confidence Sets}
  \item \hyperlink{2.8}{2.8\quad Prediction, and the Return of Conditional Expectation}
  \item \hyperlink{2.9}{2.9\quad If You Have a Prior: the Posterior Mean}
  \item \hyperlink{2.10}{2.10\quad Shrinkage and Penalization}
  \item \hyperlink{2.11}{2.11\quad Looking Ahead: Sufficiency}
\end{lecturecontents}

\bigskip
\noindent\textbf{Overview.} Estimation gets its own unit: the decision problem in which the
action space \emph{is} the parameter space and the loss is squared error. The week's deck is
an itinerary---most of its slides name what is then developed at the board---and these notes
work each named item in full. Unbiasedness is defined in its general form (an estimand
$g(\theta)$, an estimator $T(X)$) and defended as a way of choosing among admissible rules;
the information bound of \S\,\xrefL{1.9} is restated and then lifted to models with several
parameters, where the price of nuisance parameters becomes visible. The unit's centerpiece is
the course's first general-purpose \emph{constructor} of estimators: maximum likelihood,
whose estimate sits---approximately, when the information is high---at truth plus score over
information, and therefore approximately attains the floor. Four classical samples make that
exact or nearly so. Confidence sets then turn point estimates into precision statements via
pivot statistics; conditional expectation returns as the best predictor under squared error;
and when a prior is on the table, the posterior mean is \emph{proved} to minimize the Bayes
risk---the theorem whose special case the trial met in \xrefL{1.A.5}---with the
normal--normal, binomial--beta, and Poisson--gamma trios worked out. A closing pair of
slides reads Bayes methods as shrinkage and penalization, and names the next horizon:
sufficiency. \HPS{2}

\clearpage
% ============================================================
\sub{2.1}{Estimation as a Decision Problem}

Unit 1 built a general machine---action space, loss function, method, risk
(\S\,\xrefL{1.5})---and then named its two workhorse special cases (\S\,\xrefL{1.8}). This
week the first of them takes over. The deck's opening slide re-fixes the frame in three
bullets, and the first two objects below are those bullets; the third bullet---``squared
error loss is a reasonable choice''---turns out to lean on the second in a way worth making
explicit.

\begin{ObjL}{Concept}{2.1.1}{Estimation, resumed}
Estimation is the class of ``decision problems where the action space is the parameter
space'' (\xrefL{1.8.1}): the act is a guess at the unknown, and reporting $a$ when the truth
is $\theta$ costs $(a-\theta)^2$---``squared error loss is a reasonable choice.'' Everything
Unit 1 said about methods and risk applies verbatim; what is new this week is structure
special to this action space: a floor on risk, a machine for building estimators, and a way
to report precision.
\end{ObjL}

\begin{ObjL}{Remark}{2.1.2}{``Close parameters are good guesses''}
The deck's middle bullet is a condition, not a throwaway: ``parameterization satisfies:
close parameters are good guesses.'' Squared error charges by \emph{distance in the
parameter}, so it is a sensible price only when parameter distance tracks practical
distance---when being off by $0.01$ matters about equally everywhere in $\Theta$. That is a
property of the \emph{parameterization}, not of nature: the same family of distributions can
be labeled in ways that make squared error reasonable or misleading. The exponential sample
of \xref{2.5.4} will make this concrete---labeled by its mean, the model hands the obvious
estimator a clean optimality; labeled by its rate, the ``same'' estimator is not even
unbiased. Choosing the parameterization is part of posing the estimation problem.
\end{ObjL}

\begin{Obj}{Definition}{2.1.3}{Estimand and estimator}
Let $\{P_{\theta} : \theta\in\Theta\}$ be a statistical model (\xrefL{1.3.1}) and let $X$ be
a random element---a single observation, a vector, or an entire data set---whose distribution
is $P_{\theta}$ for some unknown $\theta\in\Theta$. The \emph{estimand} is a map
$g:\Theta\to\mathbb{R}^k$, the quantity to be estimated; an \emph{estimator} of $g(\theta)$
is a statistic
\[ T \;=\; T(X) , \]
a function of the data alone.\note{Letters follow the deck's typeset Definition~1: data $X$,
estimator $T$, estimand $g$ (the week-1 deck wrote data $Y$ and rules $\delta$; the Fisher
handout wrote $d(Y)$---see \xrefL{1.9.1}). The deck also switches the parameter to a
subscript, $\E_{\theta}$; these notes keep the superscript $\E^{\theta}$ of
\xrefL{1.3.1}---same object, same fixed candidate truth.}
\end{Obj}

\begin{Obj}{Definition}{2.1.4}{Unbiased estimator, general case}
The estimator $T=T(X)$ is \emph{unbiased} for $g(\theta)$ if
\[ \E^{\theta}\{T(X)\} \;=\; g(\theta) \qquad\text{for all } \theta\in\Theta ; \]
equivalently, if its \emph{bias} vanishes identically:
\[ \mathrm{Bias}^{\theta}(T) \;:=\; \E^{\theta}\{T(X)\} - g(\theta) \;=\; 0
   \qquad\text{for all } \theta\in\Theta . \]
\end{Obj}

\begin{ObjL}{Remark}{2.1.5}{Reading the definition}
The teeth are, as always, in ``for all $\theta$'': the identity must hold at every candidate
truth, not at a lucky one---this is the same quantifier that powered the differentiation
trick of \xrefL{1.9.10}. The general $g$ earns its keep at once: in the trial the estimand is
$g(\theta_T,\theta_P)=\theta_T-\theta_P$, not the parameter itself, and \xrefL{1.9.1} is
recovered by taking $g$ the identity. For $k>1$ the display is read coordinate by
coordinate. And the usual case discipline applies (\xrefL{1.1.7}): the \emph{estimator} $T$
is a random variable; the \emph{estimate} $t=T(x)$ is the number computed from your data.
\end{ObjL}

\begin{ObjL}{Example}{2.1.6}{Unbiased or not: three classics}
(a) \emph{The sample mean.} For draws $X_1,\dots,X_n$ with common expectation $\mu$,
linearity gives $\E\{\bar X\}=\tfrac1n\sum_i \E\{X_i\}=\mu$: the average is unbiased for
$\mu$ in \emph{any} model with a mean---no shape assumptions used.
(b) \emph{The observed rate.} The arm total's rate $X/n$ is unbiased for a cure probability
(\xrefL{1.A.3}), and \xrefL{1.A.6} showed it is, among unbiased estimators, unimprovable.
(c) \emph{The sample variance, and the $n-1$ mystery.} The natural plug-in
$\widehat{\sigma^2}=\frac1n\sum_i (X_i-\bar X)^2$ is \emph{biased}:
\[ \E\Big\{\tfrac1n\textstyle\sum_i (X_i-\bar X)^2\Big\} \;=\; \frac{n-1}{n}\,\sigma^2
   \;<\;\sigma^2 , \]
because deviations are measured from $\bar X$, which has chased the data part of the way
(the computation, every step, is \appref{2.A.1}). Dividing by $n-1$ instead of $n$ repairs
the average exactly---that, and nothing deeper, is where the $n-1$ comes from. The repaired
version is unbiased; whether it is \emph{better} is a risk question, not an unbiasedness
question.
\end{ObjL}

\begin{ObjL}{Concept}{2.1.7}{Why restrict to unbiased estimators?}
The deck's four bullets, in order. \emph{``Need to have a way to choose from admissible
estimators''}: admissibility (\xrefL{1.5.13}) discards the dominated but leaves a crowd whose
risk curves cross; restricting the contest is one way to get a winner. \emph{``It makes
sense''}: deadpan, but real---a rule that is systematically off at some $\theta$ is lying
about that corner of the parameter space, and forbidding that is a defensible ground rule.
\emph{``The class will typically have a unique uniformly optimum''}: within the unbiased
class, risk is variance (\xrefL{1.9.2}), and in well-behaved models one estimator has the
smallest variance at \emph{every} $\theta$ simultaneously---a \emph{uniformly} best member
(the usual name, for later, is UMVU: uniformly minimum-variance unbiased). No such uniform
winner exists among all rules (the crossing curves of \xrefL{1.5.14}); inside the class, the
contest has a champion. \emph{``Unbiased estimators in largish samples with
smallish-dimensional parameters are usually just fine''}: the honest disclaimer---with much
data and few parameters, the bias one could profitably trade away is small anyway. Unit 1's
counterpoint stands: with little data, a shaded (biased) estimator can beat every unbiased
one over most of the parameter space (\xrefL{1.A.4}). The restriction is a choice, and the
deck presents it as one.
\end{ObjL}

% ============================================================
\sub{2.2}{The Efficiency Bound, Revisited}

``We did the efficiency bound last time,'' the deck's fifth slide begins, and compresses
\S\,\xrefL{1.9} into four bullets. The compression is worth reading closely---it contains the
whole proof---and then the section does what the deck does next: it removes the assumption
that $\theta$ is the only unknown.

\begin{ObjL}{Concept}{2.2.1}{The bound in one paragraph}
The deck's recap, unpacked. ``Deviation of unbiased estimate from Score divided by
information is uncorrelated with Score (because derivative of expectation is covariance with
score)'': differentiating the unbiasedness identity showed
$\Cov^{\theta}(T,U)=1$ (\xrefL{1.9.10}, step (2)), and since $U/I(\theta)$ has
$\Cov(U/I,U)=I/I=1$ as well, the residual $T-\theta-U/I(\theta)$ has covariance
$1-1=0$ with the score: the projection geometry of \xrefL{1.9.11}. The variance of that
projection, $U/I(\theta)$, is $1/I(\theta)$---no larger than the variance of the unbiased
estimator, which carries the residual too.\note{As printed, the slide's second bullet reads
``So variance of score is less than variance of unbiased estimate''---a compression: the
variance of the \emph{score} is $I(\theta)$, typically enormous; the object meant is the
variance of the score \emph{divided by the information}, i.e.\ of the projection. The
slide's own first bullet has it right. (Recorded in the errata file.)} ``And variance of
score is the information\dots'': \xrefL{1.9.8}. And the conclusion, paraphrased (the
original carries two typos, preserved in the transcript): the variance of an unbiased
estimate is no less than the inverse information---the bound, \xrefL{1.9.10}.
\end{ObjL}

\begin{Obj}{Definition}{2.2.2}{Score vector and information matrix}
For a model $f(x;\theta)$ with $\theta=(\theta_1,\dots,\theta_p)^{\top}\in\Theta\subseteq
\mathbb{R}^p$ smooth in $\theta$, the \emph{score vector} is
\[ U(\theta) \;=\; \left( \frac{\partial}{\partial\theta_1}\log f(X;\theta),\;
   \dots,\; \frac{\partial}{\partial\theta_p}\log f(X;\theta) \right)^{\!\top} , \]
and the \emph{information matrix} is the $p\times p$ matrix
\[ I(\theta) \;=\; \E^{\theta}\big\{ U(\theta)\,U(\theta)^{\top} \big\} , \qquad
   I(\theta)_{jk} \;=\; \E^{\theta}\{ U_j(\theta)\,U_k(\theta) \} . \]
\end{Obj}

\medskip
Each component score is the one-parameter score of \xrefL{1.9.3} taken while the other
coordinates hold still, so each has mean zero by \xrefL{1.9.5} applied coordinatewise; the
information matrix is therefore the covariance matrix of the score vector, and the argument
of \xrefL{1.9.7} gives, entry by entry, the curvature form
$I(\theta)_{jk}=-\E^{\theta}\{\partial^2 \log f/\partial\theta_j\partial\theta_k\}$. The
bound now reads as follows; its proof is \xrefL{1.9.10}'s with one new ingredient, the
matrix inverse.

\begin{Obj}{Theorem}{2.2.3}{The information bound when there are other parameters}
In a regular model with $I(\theta)$ invertible, let $\hat\theta=\hat\theta(X)$ be an
unbiased estimator of the \emph{first coordinate}: $\E^{\theta}\{\hat\theta\}=\theta_1$ for
all $\theta$. Then
\[ \Var^{\theta}(\hat\theta) \;\ge\; e_1^{\top} I(\theta)^{-1} e_1
   \;=\; \big(I(\theta)^{-1}\big)_{11} , \qquad e_1=(1,0,\dots,0)^{\top} . \]
\end{Obj}
\begin{Prf}{2.2.3}{Proof of the Multiparameter Bound by Projecting onto the Span of the Score Components.}{Uses \xrefL{1.9.5}, \xrefL{1.9.10}, \xref{2.2.2}.}
\item First, one covariance per coordinate. Differentiate the unbiasedness identity
      $\int \hat\theta(x) f(x;\theta)\,dx = \theta_1$ in each $\theta_j$, passing the
      derivative under the integral and substituting
      $\partial f/\partial\theta_j = U_j f$ as in \xrefL{1.9.10}(2):
      \[ \E^{\theta}\{\hat\theta\, U_j(\theta)\} \;=\;
         \frac{\partial\,\theta_1}{\partial\theta_j} \;=\;
         \begin{cases} 1, & j=1\\ 0, & j\ne 1 \end{cases}
         \qquad\text{i.e.}\qquad
         \E^{\theta}\{\hat\theta\, U(\theta)\} \;=\; e_1 . \]
      Since each $U_j$ has mean zero, the same holds with $\hat\theta-\theta_1$ in place of
      $\hat\theta$: the deviation covaries with the score components according to $e_1$.
\item Next, the candidate projection: consider the scalar random variable
      $e_1^{\top} I(\theta)^{-1} U(\theta)$, a fixed linear combination of the score
      components. Its covariances with the components are
      \[ \E^{\theta}\big\{ \big(e_1^{\top} I^{-1} U\big)\, U^{\top} \big\}
         \;=\; e_1^{\top} I^{-1}\, \E^{\theta}\{U U^{\top}\}
         \;=\; e_1^{\top} I^{-1} I \;=\; e_1^{\top} . \]
      Identical to step (1)---so the difference is orthogonal to every score component:
      \[ \E^{\theta}\Big[ \Big( \hat\theta - \theta_1 - e_1^{\top} I(\theta)^{-1}
         U(\theta) \Big)\, U(\theta) \Big] \;=\; \mathbf{0} . \]
      This is the deck's displayed equation (its $\theta_0$ is the true parameter value, the
      estimand being its first coordinate); ``geometrically, we are projecting the unbiased
      estimator onto the span of the components of the score.''
\item The projection's variance, by squaring the linear combination:
      \[ \Var^{\theta}\big( e_1^{\top} I^{-1} U \big)
         = e_1^{\top} I^{-1}\,\E^{\theta}\{U U^{\top}\}\, I^{-1} e_1
         = e_1^{\top} I^{-1} I\, I^{-1} e_1
         = e_1^{\top} I^{-1} e_1 = \big(I^{-1}\big)_{11} . \]
\item Pythagoras finishes it: by step (2) the deviation splits into the projection plus an
      orthogonal residual, so variances add as in \xrefL{1.9.11}:
      \[ \Var^{\theta}(\hat\theta)
         \;=\; \big(I^{-1}\big)_{11} \;+\;
         \Var^{\theta}\Big( \hat\theta-\theta_1- e_1^{\top} I^{-1} U \Big)
         \;\ge\; \big(I^{-1}\big)_{11} . \]
\end{Prf}

\begin{ObjL}{Remark}{2.2.4}{The geometry, upgraded}
Nothing changed but the dimension of the stage. In \xrefL{1.9.11} the score spanned a
\emph{line} and the bound was the squared length of a projection onto it; here the $p$
score components span a $p$-dimensional subspace, the projection of $\hat\theta-\theta_1$
onto that subspace is $e_1^{\top}I^{-1}U$, and the right triangle of \xrefL{1.9.12} is
intact---hypotenuse the deviation, base the projection, and the bound again the base's
squared length.
\end{ObjL}

\begin{ObjL}{Remark}{2.2.5}{Nuisance parameters cost information}
Why $\big(I^{-1}\big)_{11}$ and not $1/I_{11}$? Because they differ, and the difference is
the price of company:
\[ \big(I^{-1}\big)_{11} \;\ge\; \frac{1}{I_{11}} ,
   \qquad\text{with equality if and only if } I_{1j}=0 \text{ for all } j\ne 1 \]
(the $2\times2$ case, with every step, is \appref{2.A.2}). If the other coordinates were
\emph{known}, the model would be one-parameter and the floor would be $1/I_{11}$; having to
carry them raises the floor unless their scores are uncorrelated with the first's---the
\emph{orthogonal} case. The normal sample of \xref{2.6.3} lands exactly there: its
information matrix is diagonal, so not knowing the variance costs the mean's estimator
nothing. When the off-diagonals do not vanish, estimating your parameter of interest is
genuinely harder because you must estimate its neighbors' influence too.
\end{ObjL}

% ============================================================
\sub{2.3}{Two Information Computations}

The deck's seventh slide names two models and then asks, flatly, ``What is information?''
The two computations first; the question, with its three answers, after.

\begin{ObjL}{Example}{2.3.1}{Binomial}
One arm of the trial, revisited with the new vocabulary: $X\sim$ Binomial$(n,\theta)$ (the
total of $n$ iid Bernoulli variables). Its information was computed, both ways and with
every step, in \xrefL{1.A.6}:
\[ I_n(\theta) \;=\; \frac{n}{\theta(1-\theta)} , \qquad
   \frac{1}{I_n(\theta)} \;=\; \frac{\theta(1-\theta)}{n} \;=\; \Var^{\theta}(X/n) : \]
the observed rate attains the floor. Note where the floor is \emph{highest}:
$\theta(1-\theta)$ peaks at $\theta=\tfrac12$, so the middle of the parameter space is the
noisiest place to estimate a proportion---and the information, its reciprocal (times $n$),
is smallest there. A coin near fair is the hardest coin to pin down.
\end{ObjL}

\begin{ObjL}{Example}{2.3.2}{Multivariate normal with known covariance}
Let $Y\sim N_p(\theta,\Sigma)$ with $\Sigma$ known and invertible: the estimand is the mean
vector itself. The log-density is
\[ \log f(y;\theta) \;=\; -\tfrac12\,(y-\theta)^{\top}\Sigma^{-1}(y-\theta)
   \;-\; \tfrac12\log\big((2\pi)^p\,|\Sigma|\big) , \]
and differentiating in $\theta$ kills the constant and unfolds the quadratic:
\[ U(\theta) \;=\; \Sigma^{-1}(Y-\theta) , \qquad
   I(\theta) \;=\; \E^{\theta}\big\{\Sigma^{-1}(Y-\theta)(Y-\theta)^{\top}\Sigma^{-1}\big\}
   \;=\; \Sigma^{-1}\,\Sigma\,\Sigma^{-1} \;=\; \Sigma^{-1} . \]
The information is the \emph{precision matrix}---the covariance's inverse---and it does not
depend on $\theta$. For an iid sample of size $n$ the information is $n\Sigma^{-1}$
(\xrefL{1.9.9}, coordinatewise), the floor for the first coordinate is
$\big((n\Sigma^{-1})^{-1}\big)_{11}=\Sigma_{11}/n$, and $\bar X_1$---with variance exactly
$\Sigma_{11}/n$---attains it.
\end{ObjL}

\begin{ObjL}{Remark}{2.3.3}{``What is information?''}
The deck's question, and the three answers the course now owns. \emph{Algebraically}: the
variance (matrix) of the score---how much the data's argument about $\theta$ swings from
sample to sample (\xrefL{1.9.8}). \emph{Geometrically}: the expected curvature of the
log-likelihood at the truth---how sharply the model distinguishes $\theta$ from its
neighbors (\xrefL{1.9.7}). \emph{Operationally}: an inverse variance---the best precision
any unbiased estimator can achieve, attained in the two examples above. And it is additive:
$n$ independent observations carry $n$ times the information of one (\xrefL{1.9.9})---the
sense in which evidence accumulates.
\end{ObjL}

% ============================================================
\sub{2.4}{Maximum Likelihood}

Everything so far judges estimators after someone proposes one; the floor of \xref{2.2.3}
says how good a proposal could possibly be, and stays silent about where proposals come
from. The deck's eighth slide supplies the factory. The idea is a reversal: the density
$f(x;\theta)$ was built to be read as a function of $x$ at a fixed $\theta$---probabilities
of data under a hypothesis. Read it the other way.

\begin{Obj}{Definition}{2.4.1}{The likelihood}
Given observed data $X$, the \emph{likelihood function} of the model $f(x;\theta)$ is
\[ \mathrm{lik}(\theta) \;=\; f(X;\theta) , \qquad \theta\in\Theta , \]
the joint density (or pmf) of the data, evaluated at the data, read as a function of the
parameter. Its logarithm $\ell(\theta)=\log \mathrm{lik}(\theta)$ is the
\emph{log-likelihood}.
\end{Obj}

\begin{ObjL}{Remark}{2.4.2}{Reading the likelihood}
$\mathrm{lik}(\theta)$ answers, for each candidate truth: \emph{how probable was what I
actually saw, if this candidate were the truth?} It ranks explanations. It is
\emph{not} a probability distribution over $\theta$---nothing says
$\int \mathrm{lik}(\theta)\,d\theta=1$, and no randomness attaches to $\theta$
here\note{The distinction is worth a sentence, because the word ``likely'' invites the
slide. ``The probability of the data given $\theta$'' is a statement about \emph{data};
the likelihood re-reads that number as a score for \emph{$\theta$}, the data now frozen. A
statement like ``$\theta$ is probably near $0.4$'' needs $\theta$ to be random---which is
Bayesian territory, where \S\,\xref{2.9}'s posterior, prior $\times$ likelihood normalized,
makes it legitimate.}---only its \emph{ratios} across candidate $\theta$ carry meaning, which
is why multiplicative constants (in $\theta$) are routinely dropped. The score of
\xrefL{1.9.3} is exactly the log-likelihood's slope: the two sections have been talking
about the same curve all along.
\end{ObjL}

\begin{Obj}{Definition}{2.4.3}{Maximum likelihood estimator}
The \emph{maximum likelihood estimator} (MLE) is the parameter value at which the
likelihood---equivalently the log-likelihood---is largest:
\[ \hat\theta \;=\; \hat\theta(X) \;=\; \argmax_{\theta\in\Theta}\, \mathrm{lik}(\theta)
   \;=\; \argmax_{\theta\in\Theta}\, \ell(\theta) . \]
\end{Obj}

\begin{ObjL}{Remark}{2.4.4}{``The score!''}
The deck's sub-bullet under ``The MLE'' is two words and an exclamation point, and it is the
right reaction. If the maximum sits in the interior of $\Theta$ and $\ell$ is smooth, the
first-order condition is
\[ 0 \;=\; \ell'(\hat\theta) \;=\; U(\hat\theta) : \]
\emph{the MLE is the parameter value at which the score vanishes}---where the data stop
arguing for a nudge in either direction (\xrefL{1.9.4}). The object built in \S\,\xrefL{1.9}
to measure estimators turns out to be the equation that constructs one.
\end{ObjL}

\begin{ObjL}{Concept}{2.4.5}{Truth plus score over information}
The deck: ``The MLE differs from parameter by the ratio of score to information
(approximately with high information).'' Here is the calculation behind the sentence. Expand
the score equation $0=\ell'(\hat\theta)$ in a Taylor series about the true value $\theta$:
\[ 0 \;=\; \ell'(\hat\theta) \;\approx\; \ell'(\theta) + (\hat\theta-\theta)\,
   \ell''(\theta) , \]
keeping the linear term. Solve for the deviation:
\[ \hat\theta - \theta \;\approx\; \frac{\ell'(\theta)}{-\ell''(\theta)}
   \;=\; \frac{U(\theta)}{-\ell''(\theta)}
   \;\approx\; \frac{U(\theta)}{I(\theta)} , \]
the last step replacing the realized curvature $-\ell''(\theta)$ by its expectation
$I(\theta)$ (\xrefL{1.9.7})---a substitution that is good when the information is high,
because the curvature, a sum of $n$ independent contributions, then concentrates near its
mean. Both approximations sharpen as $I(\theta)$ grows; neither is a theorem as stated, and
the deck's parenthetical---``approximately with high information''---is the honest label.
The notes will treat it as the deck does: a heuristic with excellent credentials, exact in
the cleanest models (\xref{2.5.1}) and provable-with-care in wide generality (a story for a
later course).
\end{ObjL}

\begin{ObjL}{Remark}{2.4.6}{What the heuristic buys}
Take the approximation $\hat\theta-\theta\approx U(\theta)/I(\theta)$ seriously and compute
with it. The right side has expectation $\E^{\theta}\{U\}/I=0$ (\xrefL{1.9.5}): the MLE is
\emph{approximately unbiased}. Its variance is
$\Var^{\theta}(U)/I(\theta)^2=I(\theta)/I(\theta)^2=1/I(\theta)$ (\xrefL{1.9.8}): the MLE
\emph{approximately attains the information bound}. That is the week's arc in one line: the
floor of \S\,\xrefL{1.9} is not just a benchmark---maximum likelihood is a general-purpose
recipe that (approximately, eventually) reaches it. And the expression $U/I$ is an old
friend: it is exactly the quantity whose \emph{exact} equality with $\hat\theta-\theta$
characterized the models where the bound is hit on the nose
(\xrefL{1.9.11}, \xrefL{1.9.13}).
\end{ObjL}

\begin{Fig}{2.4.7}{Likelihood, score, and information in one picture}{Two log-likelihoods
with the same maximizer $\hat\theta$ (dashed drop). The slope of the curve is the score,
zero at the peak (\xref{2.4.4}); at a candidate $\theta$ left of the peak the slope is
positive---the data argue for moving right. The curvature is the information's realized
version: on the left the curve is sharply bent ($-\ell''$ large), so parameter values away
from $\hat\theta$ are much worse explanations and the peak deserves trust; on the right the
same peak sits on a nearly flat ridge ($-\ell''$ small)---many candidates explain the data
almost equally well, and the estimate, while the same number, means less. The information
bound $1/I$ of \xrefL{1.9.10} is this picture's moral quantified.}
\begin{tikzpicture}[x=1cm,y=1cm, every node/.style={font=\small}]
  % ---- left: high information ----
  \node[align=center] at (2.2,2.6) {high information};
  \draw[black!60] (-0.2,0) -- (4.6,0);
  \draw[figTerm, very thick, domain=1.05:3.35, smooth, samples=60]
    plot (\x, {2.0 - 1.55*(\x-2.2)^2});
  \draw[black!55, densely dashed, thick] (2.2,0) -- (2.2,2.0);
  \node[black!70, font=\footnotesize, below] at (2.2,-0.04) {$\hat\theta$};
  % tangent segment showing positive slope (score) at theta = 1.3 < thetahat
  % curve y = 2.0 - 1.55(x-2.2)^2: point (1.3, 0.7445), slope 2.79
  \draw[black!65, thick] (1.15,0.326) -- (1.55,1.442);
  \node[black!65, font=\footnotesize, fill=white, inner sep=1.5pt, left=2pt]
    at (1.26,1.14) {slope $=U(\theta)$};
  \node[black!75, font=\footnotesize] at (2.2,-0.62) {$-\ell''$ large: peak trusted};
  % ---- right: low information ----
  \begin{scope}[shift={(6.9,0)}]
    \node[align=center] at (2.2,2.6) {low information};
    \draw[black!60] (-0.2,0) -- (4.6,0);
    \draw[figTerm, very thick, domain=0.15:4.25, smooth, samples=60]
      plot (\x, {2.0 - 0.22*(\x-2.2)^2});
    \draw[figLim, densely dashed, thick] (2.2,0) -- (2.2,2.0);
    \node[figLim!80!black, font=\footnotesize, below] at (2.2,-0.04) {$\hat\theta$};
    \node[black!75, font=\footnotesize] at (2.2,-0.62) {$-\ell''$ small: peak vague};
  \end{scope}
\end{tikzpicture}
\end{Fig}

% ============================================================
\sub{2.5}{Four Samples, Four Maximum-Likelihood Estimators}

The deck's ninth slide is four one-word prompts---``Normal, Bernoulli, Exponential,
Poisson''---and two questions: ``What is the MLE,'' and ``How does it compare to the
information inequality?'' As with Unit 1's closing checklist, the notes work what the deck
names. Each example runs the same drill: likelihood, log, score, MLE, information, and the
comparison. Throughout, $X_1,\dots,X_n$ are iid and $\bar X$ is their average.

\begin{ObjL}{Example}{2.5.1}{Normal, known variance}
$X_i\sim N(\mu,\sigma^2)$ with $\sigma^2$ known. The score was computed in \xrefL{1.9.13}:
$U(\mu)=n(\bar X-\mu)/\sigma^2$, which vanishes exactly at
\[ \hat\mu \;=\; \bar X . \]
(The log-likelihood is a downward parabola in $\mu$, so the critical point is the global
maximum.) Information $I_n(\mu)=n/\sigma^2$; floor $\sigma^2/n$; and
$\Var^{\mu}(\bar X)=\sigma^2/n$: the MLE is unbiased and attains the bound \emph{exactly, at
every} $n$---the degenerate-triangle model of \xrefL{1.9.13}, where
$\hat\mu-\mu=U(\mu)/I(\mu)$ holds with no approximation at all.
\end{ObjL}

\begin{ObjL}{Example}{2.5.2}{Bernoulli}
$X_i\sim$ Bernoulli$(\theta)$, $0<\theta<1$, with arm total $X=\sum_i X_i$. The
log-likelihood and score, from \xrefL{1.A.6}(1)--(2) summed over the sample:
\[ \ell(\theta) \;=\; X\log\theta + (n-X)\log(1-\theta) , \qquad
   U(\theta) \;=\; \frac{X}{\theta}-\frac{n-X}{1-\theta}
   \;=\; \frac{X-n\theta}{\theta(1-\theta)} . \]
The score vanishes exactly when $X=n\theta$, i.e.\ at
\[ \hat\theta \;=\; \frac{X}{n} = \bar X \]
(and $\ell''<0$ throughout, so this is the maximum; for the boundary samples $X=0$ and
$X=n$ the likelihood is monotone, and its maximum sits at the same $X/n$, now an endpoint).
The observed rate---Unit 1's estimator of choice---\emph{is} the MLE. Information $I_n(\theta)=n/(\theta(1-\theta))$
(\xrefL{1.A.6}); floor $\theta(1-\theta)/n$; variance of the rate: the same
(\xrefL{1.A.3}). Attained exactly.
\end{ObjL}

\begin{ObjL}{Example}{2.5.3}{Poisson}
$X_i\sim$ Poisson$(\lambda)$, $\lambda>0$: pmf $f(x;\lambda)=e^{-\lambda}\lambda^x/x!$, so
\[ \ell(\lambda) \;=\; -n\lambda + \Big(\sum_i X_i\Big)\log\lambda - \sum_i \log(X_i!) ,
   \qquad
   U(\lambda) \;=\; -n + \frac{\sum_i X_i}{\lambda}
   \;=\; \frac{\sum_i X_i - n\lambda}{\lambda} . \]
For samples with $\sum_i X_i>0$ the score vanishes at $\hat\lambda=\bar X$, and the second
derivative $-\sum_i X_i/\lambda^2<0$ confirms a maximum (an all-zero sample makes the
likelihood $e^{-n\lambda}$ monotone decreasing, driving the maximizer to the boundary,
where $\bar X=0$ still names it). Information, by the curvature form
(\xrefL{1.9.7}): $I_n(\lambda)=\E\{\sum_i X_i\}/\lambda^2=n\lambda/\lambda^2=n/\lambda$.
Floor $\lambda/n$; and since a Poisson variable has variance equal to its mean,
$\Var(\bar X)=\lambda/n$. Attained exactly---and indeed
$U(\lambda)/I_n(\lambda)=\frac{\sum_i X_i-n\lambda}{\lambda}\cdot\frac{\lambda}{n}
=\bar X-\lambda$, the equality case once more.
\end{ObjL}

\begin{ObjL}{Example}{2.5.4}{Exponential: one sample, two parameterizations}
Exponential variables can be labeled by their \emph{mean} $\mu$ (density
$f(x;\mu)=\mu^{-1}e^{-x/\mu}$, $x>0$) or by their \emph{rate} $\lambda=1/\mu$ (density
$\lambda e^{-\lambda x}$). The choice, innocuous-looking, decides how estimation goes---the
promise of \xref{2.1.2} kept.

\emph{Mean parameterization.}
\[ \ell(\mu) \;=\; -n\log\mu - \frac{n\bar X}{\mu} , \qquad
   U(\mu) \;=\; -\frac{n}{\mu} + \frac{n\bar X}{\mu^2}
   \;=\; \frac{n(\bar X-\mu)}{\mu^2} , \]
so $\hat\mu=\bar X$: unbiased (linearity). Information: differentiating,
$\ell''(\mu)=n/\mu^2-2n\bar X/\mu^3$, so
$I_n(\mu)=-\E\{\ell''(\mu)\}=-n/\mu^2+2n\mu/\mu^3=n/\mu^2$; the floor is $\mu^2/n$; and
since an exponential's variance is its squared mean, $\Var(\bar X)=\mu^2/n$. Attained
exactly ($U/I_n=\bar X-\mu$ once again).

\emph{Rate parameterization.}
\[ \ell(\lambda) \;=\; n\log\lambda - \lambda\, n\bar X , \qquad
   U(\lambda) \;=\; \frac{n}{\lambda} - n\bar X , \]
so $\hat\lambda=1/\bar X$---by invariance the natural counterpart, but \emph{not} unbiased:
\[ \E^{\lambda}\{\hat\lambda\} \;=\; \frac{n\lambda}{n-1} \;>\; \lambda
   \qquad (n\ge 2) , \]
computed through the Gamma distribution of the sample total in \appref{2.A.3}. The bound of
\xrefL{1.9.10} does not even apply to it as stated; and the repaired, unbiased version
$(n-1)/(n\bar X)$ has variance $\lambda^2/(n-2)$ (\appref{2.A.3})---strictly \emph{above}
the floor $\lambda^2/n$, at every $n\ge3$. Same data, same principle; the parameterization
alone separates ``exactly optimal'' from ``biased, and above the floor even after repair.''
\end{ObjL}

\begin{ObjL}{Remark}{2.5.5}{The pattern, and its limits}
The deck's second question---``How does it compare to the information inequality?''---now
has four answers:
\begin{center}\footnotesize
\begin{tabular}{llccl}
\hline
model & MLE & $I_n$ & floor $1/I_n$ & bound attained? \\
\hline
Normal mean ($\sigma^2$ known) & $\bar X$ & $n/\sigma^2$ & $\sigma^2/n$ & yes---exactly \\
Bernoulli $\theta$ & $X/n$ & $n/(\theta(1-\theta))$ & $\theta(1-\theta)/n$ & yes---exactly \\
Poisson $\lambda$ & $\bar X$ & $n/\lambda$ & $\lambda/n$ & yes---exactly \\
Exponential mean $\mu$ & $\bar X$ & $n/\mu^2$ & $\mu^2/n$ & yes---exactly \\
Exponential rate $\lambda$ & $1/\bar X$ & $n/\lambda^2$ & $\lambda^2/n$ &
  no---biased; repaired, still above \\
\hline
\end{tabular}
\end{center}
\smallskip
The run of ``exactly'' is not the general situation---it is the luck of these particular
families under these particular labelings, the models where $\theta+U(\theta)/I(\theta)$
happens to assemble into a statistic (\xrefL{1.9.11})---here, each time, into $\bar X$
itself. The exponential-rate row shows how thin the luck is: a
smooth relabeling of the \emph{same} model breaks exactness. What survives relabeling is
\xref{2.4.6}'s asymptotic version---approximately unbiased, approximately at the floor, when
the information is high---and that, not exactness, is maximum likelihood's true general
credential.
\end{ObjL}

% ============================================================
\sub{2.6}{Several Parameters at Once}

Real models rarely have a single unknown. The deck's tenth slide asks for the
multi-parameter version of \xref{2.4.5} and of its variance consequence, then names two
examples; \xref{2.2.2}'s score vector and information matrix are already on the table.

\begin{ObjL}{Concept}{2.6.1}{Taylor series expansion for the MLE}
For $\theta\in\Theta\subseteq\mathbb{R}^p$, the MLE solves the $p$ score equations
$U(\hat\theta)=\mathbf 0$. Expanding each coordinate of $U$ about the truth, to first
order,
\[ \mathbf 0 \;=\; U(\hat\theta) \;\approx\; U(\theta) +
   \nabla U(\theta)\,(\hat\theta-\theta) , \]
where $\nabla U(\theta)$ is the $p\times p$ matrix of second partials of $\ell$. Replacing
that realized curvature matrix by its expectation $-I(\theta)$ (\xref{2.2.2}) and solving,
\[ \hat\theta - \theta \;\approx\; I(\theta)^{-1}\, U(\theta) : \]
truth plus \emph{matrix} inverse of information times score---\xref{2.4.5} with the division
promoted to a solve. The same two approximations, with the same ``high information'' license.
\end{ObjL}

\begin{ObjL}{Concept}{2.6.2}{Covariance matrix of the multivariate MLE}
Compute with the approximation. The score vector has mean $\mathbf 0$ and covariance matrix
$I(\theta)$ (\xref{2.2.2}), so $\hat\theta$ is approximately unbiased and
\[ \Cov(\hat\theta) \;\approx\; I^{-1}\,\E\{U U^{\top}\}\, I^{-1}
   \;=\; I^{-1}\, I\, I^{-1} \;=\; I(\theta)^{-1} : \]
the full inverse information matrix. In particular
$\Var(\hat\theta_1)\approx (I^{-1})_{11}$---each coordinate of the MLE approximately attains
the several-parameter floor of \xref{2.2.3}, nuisance tax included.
\end{ObjL}

\begin{ObjL}{Example}{2.6.3}{Normal sample, expectation and variance both unknown}
$X_i\sim N(\mu,\sigma^2)$, both parameters unknown: $\theta=(\mu,\sigma^2)$. The
log-likelihood is
\[ \ell(\mu,\sigma^2) \;=\; -\frac{n}{2}\log(2\pi\sigma^2)
   -\frac{1}{2\sigma^2}\sum_{i=1}^n (X_i-\mu)^2 , \]
with the two score components
\[ \frac{\partial\ell}{\partial\mu} \;=\; \frac{1}{\sigma^2}\sum_i (X_i-\mu) , \qquad
   \frac{\partial\ell}{\partial(\sigma^2)} \;=\; -\frac{n}{2\sigma^2}
   +\frac{1}{2\sigma^4}\sum_i (X_i-\mu)^2 . \]
Setting both to zero: the first gives $\hat\mu=\bar X$ regardless of $\sigma^2$; substituting
into the second and solving, $\frac{n}{2\sigma^2}=\frac{1}{2\sigma^4}\sum_i(X_i-\bar X)^2$,
so
\[ \hat\mu \;=\; \bar X , \qquad
   \widehat{\sigma^2} \;=\; \frac1n\sum_{i=1}^n (X_i-\bar X)^2 . \]
Note which variance estimator maximum likelihood picked: the \emph{biased} $1/n$ version of
\xref{2.1.6}(c)---the two principles of this unit genuinely disagree here, by the factor
$(n-1)/n$. For the information matrix, differentiate the scores once more and take
expectations (using $\E\{X_i-\mu\}=0$ and $\E\{(X_i-\mu)^2\}=\sigma^2$):
\[ -\E\!\left[\frac{\partial^2 \ell}{\partial\mu^2}\right] = \frac{n}{\sigma^2}, \qquad
   -\E\!\left[\frac{\partial^2 \ell}{\partial\mu\,\partial(\sigma^2)}\right]
   = \E\!\left[\frac{\sum_i(X_i-\mu)}{\sigma^4}\right] = 0 , \qquad
   -\E\!\left[\frac{\partial^2 \ell}{\partial(\sigma^2)^2}\right]
   = -\frac{n}{2\sigma^4}+\frac{n\sigma^2}{\sigma^6} = \frac{n}{2\sigma^4} , \]
so
\[ I(\mu,\sigma^2) \;=\;
   \begin{pmatrix} \dfrac{n}{\sigma^2} & 0\\[6pt] 0 & \dfrac{n}{2\sigma^4} \end{pmatrix} ,
   \qquad
   I(\mu,\sigma^2)^{-1} \;=\;
   \begin{pmatrix} \dfrac{\sigma^2}{n} & 0\\[6pt] 0 & \dfrac{2\sigma^4}{n} \end{pmatrix} . \]
The matrix is \emph{diagonal}: the two scores are uncorrelated, the parameters are
orthogonal, and $(I^{-1})_{11}=\sigma^2/n=1/I_{11}$---by \xref{2.2.5}, not knowing the
variance costs the mean's estimator nothing (to this order). The floor for $\hat\mu$ is the
same $\sigma^2/n$ as when $\sigma^2$ was known, and $\bar X$ still attains it.
\end{ObjL}

\begin{ObjL}{Example}{2.6.4}{Multivariate normal with known covariance, again}
For an iid sample from $N_p(\theta,\Sigma)$, \xref{2.3.2} gives $I_n(\theta)=n\Sigma^{-1}$;
the MLE solves $\sum_i \Sigma^{-1}(X_i-\theta)=\mathbf 0$, i.e.\ $\hat\theta=\bar X$, whose
covariance matrix is exactly $\Sigma/n=(n\Sigma^{-1})^{-1}=I_n(\theta)^{-1}$. Here
\xref{2.6.2}'s approximation is an identity---the multivariate analogue of the degenerate
triangle.
\end{ObjL}

% ============================================================
\sub{2.7}{Confidence Sets}

An estimate alone is a bare number: $\hat\theta=0.31$ does not say whether the data would
have supported $0.30$ and $0.32$ almost equally (low information) or ruled them out (high).
The deck's eleventh slide asks for the definition and for its reading ``as a measure of
precision''; the twelfth builds the standard construction; the thirteenth, a title with no
body, names a duality the testing unit will cash in.

\begin{Obj}{Definition}{2.7.1}{Confidence set}
A \emph{confidence set} of level $1-\alpha$ for the parameter $\theta$ is a set $S(X)$,
computed from the data alone, such that
\[ P^{\theta}\big\{\, \theta \in S(X) \,\big\} \;\ge\; 1-\alpha
   \qquad\text{for all } \theta\in\Theta . \]
When $S(X)$ is an interval it is called a \emph{confidence interval}.
\end{Obj}

\begin{ObjL}{Remark}{2.7.2}{Reading ``confidence''}
Look hard at the display: the random object inside the braces is $S(X)$, \emph{not}
$\theta$. The parameter sits still; the \emph{set} jumps around with the data; the
guarantee---at every candidate truth, and before the data arrive---is that the procedure's
set captures the truth with probability at least $1-\alpha$. It is a statement about the
\emph{method}, in exactly \xrefL{1.5.6}'s sense of a plan graded before the data
arrive.\note{Hence the classical caution: after computing, say, the interval $[0.22,0.40]$,
the statement ``$\theta$ is in $[0.22,0.40]$ with probability $0.95$'' is not licensed---%
nothing random remains in it. The honest post-data sentence is operational: the interval was
produced by a procedure that traps the truth $95\%$ of the time. A probability distribution
\emph{for $\theta$ itself} requires a prior, and then \S\,\xref{2.9}'s posterior supplies
genuine probability statements of just that kind.}
And the deck's phrase ``as a measure of precision'' is the point of reporting the whole set:
its \emph{size} is the resolution of the study. \xref{2.7.5} will make that literal---the
width shrinks exactly like the square root of the inverse information.
\end{ObjL}

\begin{Obj}{Definition}{2.7.3}{Pivot statistic}
A \emph{pivot} is a function $T(X,\theta)$ of the data \emph{and} the parameter whose
distribution under $P^{\theta}$ is the same known distribution for every
$\theta\in\Theta$.
\end{Obj}

\begin{ObjL}{Remark}{2.7.4}{From pivot to confidence set}
A pivot converts distributional knowledge into coverage. Choose a region $A$ with
$P\{T(X,\theta)\in A\}=1-\alpha$---possible because the distribution of the pivot is known
and free of $\theta$---and \emph{invert}: define
\[ S(X) \;=\; \{\, \theta : T(X,\theta)\in A \,\} . \]
Then for every $\theta$,
\[ P^{\theta}\{\theta\in S(X)\} \;=\; P^{\theta}\{T(X,\theta)\in A\} \;=\; 1-\alpha : \]
the two events are the \emph{same} event, described from the two ends of the
$(\text{data},\text{parameter})$ pair. All the standard intervals are this one identity
wearing different pivots.
\end{ObjL}

\begin{ObjL}{Example}{2.7.5}{The $z$-interval: normal mean, known variance}
$X_i\sim N(\mu,\sigma^2)$, $\sigma^2$ known. Since $\bar X\sim N(\mu,\sigma^2/n)$,
standardizing gives
\[ Z \;=\; \frac{\sqrt n\,(\bar X-\mu)}{\sigma} \;\sim\; N(0,1)
   \qquad\text{under } P^{\mu},\ \text{for every } \mu : \]
a pivot. Take $A=[-z_{\alpha/2},z_{\alpha/2}]$, the central interval catching probability
$1-\alpha$ of the standard normal, and invert the two inequalities in $\mu$:
\[ -z_{\alpha/2} \le \frac{\sqrt n\,(\bar X-\mu)}{\sigma} \le z_{\alpha/2}
   \quad\Longleftrightarrow\quad
   \mu \;\in\; \Big[\, \bar X - z_{\alpha/2}\tfrac{\sigma}{\sqrt n},\;
   \bar X + z_{\alpha/2}\tfrac{\sigma}{\sqrt n} \,\Big] . \]
The width is $2z_{\alpha/2}\,\sigma/\sqrt n = 2z_{\alpha/2}\sqrt{1/I_n(\mu)}$: the
information of \S\,\xrefL{1.9} is not only the variance floor but the \emph{ruler} for the
precision report. Quadrupling the data halves the width.
\end{ObjL}

\begin{ObjL}{Example}{2.7.6}{The $t$ statistic: normal mean, unknown variance}
With $\sigma^2$ unknown, the $z$ pivot is not a statistic-and-parameter pair we can use---%
$\sigma$ is not known. Replace it by the sample's own scale,
$s^2=\frac{1}{n-1}\sum_i(X_i-\bar X)^2$ (the unbiased version of \xref{2.1.6}(c)). The
resulting ratio
\[ T \;=\; \frac{\sqrt n\,(\bar X-\mu)}{s} \]
is again a pivot: its distribution---Student's $t$ with $n-1$ degrees of freedom---is the
same for every $(\mu,\sigma^2)$ (the classical distribution theory of the normal sample is
taken from the text here rather than re-derived, \HPS{2}). Inverting exactly as before,
\[ \mu \;\in\; \Big[\, \bar X - t_{n-1,\alpha/2}\tfrac{s}{\sqrt n},\;
   \bar X + t_{n-1,\alpha/2}\tfrac{s}{\sqrt n} \,\Big] , \]
with $t_{n-1,\alpha/2}>z_{\alpha/2}$: the interval is systematically wider than the known-%
$\sigma$ interval. That widening is \xref{2.2.5}'s nuisance tax, paid in quantiles---not
knowing $\sigma^2$ costs nothing in the information matrix (which is diagonal,
\xref{2.6.3}) but does cost in having to estimate the ruler itself; the surcharge
$t_{n-1,\alpha/2}-z_{\alpha/2}$ fades as $n$ grows.
\end{ObjL}

\begin{ObjL}{Example}{2.7.7}{Confidence sets for ratios}
The deck's last worked prompt, and the reason it says confidence \emph{sets}. Two
independent sample means estimate $\mu_1$ and $\mu_2$, say $\bar X\sim N(\mu_1,\sigma^2/n)$
and $\bar Y\sim N(\mu_2,\sigma^2/m)$ with $\sigma^2$ known, and the estimand is the ratio
$\rho=\mu_1/\mu_2$, with $\mu_2\ne0$ so that the ratio is defined (a treatment-to-control
ratio, a slope through the origin). For each
candidate $\rho$, under any truth with $\mu_1=\rho\,\mu_2$ the combination
$\bar X-\rho\,\bar Y$ is normal with mean zero and known variance
$\sigma^2(1/n+\rho^2/m)$, so
\[ Z(\rho) \;=\; \frac{\bar X-\rho\,\bar Y}{\sigma\sqrt{1/n+\rho^2/m}} \;\sim\; N(0,1) , \]
and the inversion recipe of \xref{2.7.4} yields
\[ S(X,Y) \;=\; \Big\{\, \rho :\ (\bar X-\rho\,\bar Y)^2 \;\le\;
   z_{\alpha/2}^2\,\sigma^2\big(\tfrac1n+\tfrac{\rho^2}{m}\big) \Big\} . \]
The condition is a \emph{quadratic} inequality in $\rho$ (Fieller's construction), and the
shape of its solution set depends on the leading coefficient
$\bar Y^2-z_{\alpha/2}^2\sigma^2/m$: an interval when $\bar Y$ is convincingly nonzero; two
half-lines, or all of $\mathbb{R}$, when it is not (\appref{2.A.5} runs the cases). The
strange shapes are honesty, not pathology: when the denominator's sign is in doubt, the data
genuinely cannot bound the ratio, and a procedure with guaranteed coverage must say
so.\note{An interval-valued report would have to fail its coverage promise for some
$(\mu_1,\mu_2)$; the exotic sets are the price of the ``for all $\theta$'' in
\xref{2.7.1}.}
\end{ObjL}

\begin{ObjL}{Concept}{2.7.8}{Duality with families of point-null tests}
The deck's thirteenth slide is a bare title: ``Duality between families of point null tests
and confidence sets.'' The one-sentence preview: deciding whether each candidate value
$\theta_0$ is \emph{plausible} (a test of the point null $\theta=\theta_0$, in the sense of
\xrefL{1.8.2}) and collecting the plausible ones is the same act as reporting a confidence
set---$\theta_0\in S(X)$ exactly when the test of $\theta_0$ accepts. \xref{2.7.7} was
secretly this: a family of $z$-tests, one per candidate $\rho$, inverted. The testing unit
will make the correspondence precise and run it in both directions.
\end{ObjL}

\begin{Fig}{2.7.9}{Coverage is a property of the procedure}{Fourteen replications of the
same experiment, each yielding an interval (dot at the estimate). The truth $\theta$ (dashed
green) never moves; the intervals do. Twelve of the fourteen catch $\theta$; two (red) miss.
``$95\%$ confidence'' is a statement about the long-run frequency of catches over
replications---not about any single computed interval, which either contains $\theta$ or
does not (\xref{2.7.2}).}
\begin{tikzpicture}[x=1cm,y=1cm, every node/.style={font=\small}]
  \draw[figLim, densely dashed, thick] (3,-0.15) -- (3,6.75);
  \node[figLim!80!black, font=\footnotesize, above] at (3,6.75) {$\theta$};
  % covering intervals (blue)
  \foreach \a/\b/\y in {1.8/3.6/0.50, 2.3/4.3/0.95, 2.1/3.7/1.40, 2.0/4.2/2.30,
                        1.65/3.35/2.75, 2.5/4.3/3.20, 1.75/3.85/3.65, 2.35/3.75/4.10,
                        1.65/3.55/4.55, 2.4/4.0/5.45, 1.95/3.95/5.90, 2.3/4.4/6.35}{
    \draw[figTerm, thick] (\a,\y) -- (\b,\y);
    \fill[figTerm] ({(\a+\b)/2},\y) circle (1.5pt);
  }
  % missing intervals (red)
  \foreach \a/\b/\y in {3.1/5.0/1.85, 0.85/2.85/5.00}{
    \draw[figBad, thick] (\a,\y) -- (\b,\y);
    \fill[figBad] ({(\a+\b)/2},\y) circle (1.5pt);
  }
\end{tikzpicture}
\end{Fig}

% ============================================================
\sub{2.8}{Prediction, and the Return of Conditional Expectation}

``More probability review,'' the fourteenth slide announces, modestly. Its two questions are
the engine of everything after: ``How to minimize the expected prediction error under
squared error loss?''---and the same again ``when you have a predictor?'' The answers are
two small theorems about squared error, and the second is conditional expectation's second
star turn (after \xrefL{1.4.16}).

\begin{Obj}{Theorem}{2.8.1}{The best constant predictor is the mean}
Let $Y$ have $\E\{Y^2\}<\infty$. Among all constants $c$,
\[ \E\big\{(Y-c)^2\big\} \;=\; \Var(Y) + \big(\E\{Y\}-c\big)^2 \]
is minimized uniquely at $c=\E\{Y\}$, with minimum value $\Var(Y)$.
\end{Obj}
\begin{Prf}{2.8.1}{Proof of the Best Constant by Adding and Subtracting the Mean.}{Uses \xrefL{1.A.1}.}
\item Route the error through the mean: $Y-c=(Y-\E\{Y\})+(\E\{Y\}-c)$.
\item Square the two-term sum:
      \[ (Y-c)^2 \;=\; (Y-\E\{Y\})^2 + 2\,(Y-\E\{Y\})(\E\{Y\}-c) + (\E\{Y\}-c)^2 . \]
\item Take expectations term by term. The middle term dies---$\E\{Y\}-c$ is a constant and
      deviations about the mean average to zero---leaving the display:
      $\E\{(Y-c)^2\}=\Var(Y)+(\E\{Y\}-c)^2$.
\item The first term does not involve $c$; the second is a square, zero exactly at
      $c=\E\{Y\}$ and positive elsewhere.
\end{Prf}

\begin{Obj}{Theorem}{2.8.2}{The best predictor using $X$ is the conditional expectation}
Let $(X,Y)$ be jointly distributed with $\E\{Y^2\}<\infty$. Among all functions $h$ of the
predictor $X$,
\[ \E\big\{ \big(Y-h(X)\big)^2 \big\} \]
is minimized by
\[ h(X) \;=\; \E\{Y\mid X\} . \]
\end{Obj}
\begin{Prf}{2.8.2}{Proof of the Best Predictor by Conditioning, Slice by Slice.}{Uses \xrefL{1.4.14}, \xrefL{1.4.16}, \xref{2.8.1}.}
\item Iterate the expectation through the predictor (\xrefL{1.4.16}):
      \[ \E\big\{(Y-h(X))^2\big\}
         \;=\; \E\Big[\, \E\big\{ (Y-h(X))^2 \bigm| X \big\} \,\Big] . \]
\item Fix a slice $X=x$. Inside it, $h(x)$ is a \emph{constant}, and the conditional
      distribution of $Y$ given $X=x$ is a genuine probability distribution
      (\xrefL{1.4.14})---so \xref{2.8.1} applies verbatim \emph{within the slice}:
      \[ \E\big\{ (Y-h(x))^2 \bigm| X=x \big\}
         \;\ge\; \Var(Y\mid X=x) , \]
      with equality exactly at $h(x)=\E\{Y\mid X=x\}$.
\item The choice $h(x)=\E\{Y\mid X=x\}$ therefore minimizes the inner expectation at
      \emph{every} $x$ simultaneously; averaging over the slices (step (1)) preserves the
      domination, so no other $h$ can do better overall.
\end{Prf}

\begin{ObjL}{Remark}{2.8.3}{The regression function}
The minimizer $x\mapsto \E\{Y\mid X=x\}$ has a name---the \emph{regression function} of $Y$
on $X$---and the two theorems assign the season's roles: with no predictor in hand, report
the mean; with one, report the conditional mean. This is \emph{prediction} (the target $Y$
is a random quantity not yet seen), a cousin of estimation (the target $g(\theta)$ is a
fixed unknown); the machinery of the regression weeks will be theorem \xref{2.8.2} worked
into linear form. More immediately: the deck placed this ``review'' exactly where it did
because the Bayes section next door is \xref{2.8.2} in disguise.
\end{ObjL}

% ============================================================
\sub{2.9}{If You Have a Prior: the Posterior Mean}

``What if you have a prior?'' Unit 1 built the Bayesian scaffolding---prior, Bayes risk,
Bayes principle (\S\,\xrefL{1.6})---and, in the trial, found by direct computation that the
Bayes rule was a posterior mean (\xrefL{1.A.5}). The deck now asks the structural question
outright: ``How do we minimize the Bayes risk with squared error loss? (The posterior
expectation!)''---and, just above it, plants the hint: ``Is this the same structure?'' It
is. In the Bayesian computation the parameter is a random variable, the data are the
predictor, and \S\,\xref{2.8} has already done all the work.

\begin{Obj}{Definition}{2.9.1}{The posterior distribution}
Let the parameter have prior density $\pi(\theta)$ (\xrefL{1.6.1}) and the data have
conditional density $f(x\mid\theta)$ given $\theta$.\note{Notation shift, deliberate: in
Bayesian computations $\theta$ is a random variable, so the model density $f(x;\theta)$ is
now honestly a \emph{conditional} density and is written $f(x\mid\theta)$---the same
numbers, upgraded punctuation (cf.\ \xrefL{1.6.2}).} The \emph{posterior distribution} of
$\theta$ given the observed data $X=x$ is the conditional distribution of $\theta$ given
$X=x$, with density
\[ \pi(\theta\mid x) \;=\;
   \frac{\pi(\theta)\, f(x\mid\theta)}{\displaystyle\int_{\Theta}
   \pi(t)\, f(x\mid t)\,dt} . \]
\end{Obj}

\begin{ObjL}{Remark}{2.9.2}{Reading the posterior}
The numerator is \emph{prior times likelihood}: the data enter the update through
$\mathrm{lik}(\theta)=f(x\mid\theta)$ and through nothing else---\xref{2.4.1}'s object turns
out to be the entire interface between data and belief. The denominator does not involve
$\theta$: it is one number, the normalizer that makes the slice a probability distribution
(\xrefL{1.4.14} again), which is why working Bayesians write
$\pi(\theta\mid x)\propto\pi(\theta)\,f(x\mid\theta)$ and match kernels, as every example
below will. And the posterior is the honest owner of sentences the likelihood could not
license (\xref{2.4.2}): given a prior, ``$\theta$ is probably near $0.4$'' is now a
statement with a probability attached---computed from $\pi(\cdot\mid x)$.
\end{ObjL}

\begin{Obj}{Theorem}{2.9.3}{The posterior mean minimizes the Bayes risk}
Under squared-error loss, and provided the prior gives $\theta$ a finite second moment
($\E\{\theta^2\}<\infty$), among all estimators $\delta(X)$ the Bayes risk (\xrefL{1.6.2})
\[ r(\delta) \;=\; \E\big\{ \big(\delta(X)-\theta\big)^2 \big\}
   \qquad\text{($\theta$ and $X$ jointly distributed: prior, then model)} \]
is minimized by the \emph{posterior mean},
\[ \delta(x) \;=\; \E\{\theta\mid X=x\} . \]
\end{Obj}
\begin{Prf}{2.9.3}{Proof of the Posterior-Mean Rule by Recognizing the Same Structure.}{Uses \xrefL{1.4.16}, \xref{2.8.1}, \xref{2.8.2}.}
\item Under the joint distribution of $(\theta,X)$---prior for $\theta$, model for $X$ given
      $\theta$---the Bayes risk of a rule $\delta$ is exactly a mean squared
      \emph{prediction} error: the quantity being guessed, $\theta$, is a random variable,
      and the guess, $\delta(X)$, is a function of the observed one. This is the deck's
      ``same structure'': \xref{2.8.2} with $(X,Y)$ played by $(X,\theta)$.
\item Apply \xref{2.8.2}'s argument. Iterate the expectation through the data
      (\xrefL{1.4.16}):
      \[ r(\delta) \;=\; \E\Big[\, \E\big\{ (\delta(X)-\theta)^2 \bigm| X \big\} \,\Big] . \]
\item Fix the slice $X=x$: inside it $\delta(x)$ is a constant, and the conditional
      distribution of $\theta$ given $X=x$ is the \emph{posterior} (\xref{2.9.1}). By
      \xref{2.8.1} applied to that distribution, the inner expectation
      \[ \E\big\{ (\theta-\delta(x))^2 \bigm| X=x \big\}
         \;=\; \Var(\theta\mid X=x) + \big(\E\{\theta\mid X=x\}-\delta(x)\big)^2 \]
      is minimized, slice by slice, at $\delta(x)=\E\{\theta\mid X=x\}$.
\item Minimizing inside every slice minimizes the average over slices, so the posterior mean
      minimizes $r$---the argument that ran once, in miniature and for the uniform prior, in
      \xrefL{1.A.5}(5).
\end{Prf}

\medskip
So the Bayes estimate is: \emph{update, then average}. What remains is to see the update in
closed form, and the deck names the three classical cases---``Normal, normal; Binomial,
Beta; Poisson, Gamma''---each a \emph{conjugate} pair, in which the posterior lands back in
the prior's own family and the posterior mean is a weighted compromise.

\begin{ObjL}{Example}{2.9.4}{Normal--normal}
Prior $\theta\sim N(m,\tau^2)$; one observation $X\mid\theta\sim N(\theta,\sigma^2)$ with
$\sigma^2$ known (a sample of size $n$ enters as $\bar X$ with $\sigma^2/n$ in place of
$\sigma^2$). Matching kernels---the full completing-the-square, every step, is
\appref{2.A.4}---the posterior is again normal,
\[ \theta \mid X=x \;\sim\; N\big( m^{*},\, v \big), \qquad
   \frac{1}{v} \;=\; \frac{1}{\tau^2}+\frac{1}{\sigma^2} , \qquad
   m^{*} \;=\; v\left( \frac{m}{\tau^2} + \frac{x}{\sigma^2} \right)
   \;=\; \frac{\sigma^2\, m + \tau^2\, x}{\sigma^2+\tau^2} . \]
Read it in \emph{precisions} (inverse variances, the information's units,
\xrefL{1.9.8}): precisions add, and the posterior mean weighs prior guess and data by their
precisions. Equivalently, the Bayes estimate
\[ \E\{\theta\mid X=x\} \;=\; m + \frac{\tau^2}{\sigma^2+\tau^2}\,(x-m) \]
moves from the prior mean toward the data, covering only the fraction
$\tau^2/(\sigma^2+\tau^2)$ of the distance: it \emph{shrinks} the observation toward $m$.
Vague prior ($\tau^2\to\infty$): the estimate is $x$, the data win. Precise prior
($\tau^2\to0$): the estimate is $m$, the data are ignored. Everything in between is a
negotiation at rates set by the two variances.
\end{ObjL}

\begin{ObjL}{Example}{2.9.5}{Binomial--beta}
Prior $\theta\sim$ Beta$(\alpha,\beta)$ (\xrefL{1.8.5}); data $X\mid\theta\sim$
Binomial$(n,\theta)$. Match kernels in $\theta$:
\[ \pi(\theta\mid x) \;\propto\;
   \underbrace{\theta^{\alpha-1}(1-\theta)^{\beta-1}}_{\text{prior}}\;\cdot\;
   \underbrace{\theta^{x}(1-\theta)^{n-x}}_{\text{likelihood}}
   \;=\; \theta^{(\alpha+x)-1}\,(1-\theta)^{(\beta+n-x)-1} , \]
the Beta$(\alpha+x,\;\beta+n-x)$ kernel: successes reinforce $\alpha$, failures reinforce
$\beta$. By the Beta mean formula (\xrefL{1.8.5}), the Bayes estimate is
\[ \E\{\theta\mid X=x\} \;=\; \frac{\alpha+x}{\alpha+\beta+n}
   \;=\; \frac{(\alpha+\beta)\cdot\dfrac{\alpha}{\alpha+\beta} \;+\; n\cdot\dfrac{x}{n}}
   {(\alpha+\beta)+n} : \]
a weighted average of the prior mean $\alpha/(\alpha+\beta)$ and the observed rate $x/n$,
with weights $\alpha+\beta$ and $n$---\emph{the prior is worth $\alpha+\beta$
observations}. And with $\alpha=\beta=1$ (the uniform prior), the estimate is
$(1+x)/(2+n)$: Unit 1's shaded estimator, whose origin \xrefL{1.A.5}(4) computed ad hoc, now
a two-line corollary of the general theorem---the promise ``where $(1+X)/(2+n)$ comes
from,'' kept and generalized.
\end{ObjL}

\begin{ObjL}{Example}{2.9.6}{Poisson--gamma}
Prior $\theta\sim$ Gamma$(a,b)$, density proportional to $\theta^{a-1}e^{-b\theta}$ on
$(0,\infty)$ (mean $a/b$); data $X_1,\dots,X_n\mid\theta$ iid Poisson$(\theta)$, likelihood
proportional to $\theta^{\sum_i x_i} e^{-n\theta}$. The product is the
Gamma$\big(a+\sum_i x_i,\; b+n\big)$ kernel, so
\[ \E\Big\{\theta \Bigm| \textstyle\sum_i X_i = s\Big\} \;=\; \frac{a+s}{b+n}
   \;=\; \frac{b\cdot\dfrac{a}{b} + n\cdot\bar x}{b+n} : \]
prior mean weighted $b$, sample mean weighted $n$. The prior acts like $b$ phantom
observations totaling $a$ events.
\end{ObjL}

\begin{ObjL}{Remark}{2.9.7}{One anatomy, three posteriors}
All three estimates are the same sentence:
\[ \text{posterior mean} \;=\;
   \frac{(\text{prior worth})\times(\text{prior guess})
   \;+\; (\text{data worth})\times(\text{data guess})}
   {(\text{prior worth}) + (\text{data worth})} , \]
with ``worth'' measured in observation-equivalents (or precisions, \xref{2.9.4}). As data
accumulate the data's worth grows without bound and the Bayes estimate slides toward the
data's guess; with little data the prior holds it back. That systematic holding-back has a
name, and the deck's next slide says it.
\end{ObjL}

% ============================================================
\sub{2.10}{Shrinkage and Penalization}

``By the way,'' the seventeenth slide begins---and then quietly names the modern half of the
subject.

\begin{ObjL}{Concept}{2.10.1}{Bayes as penalized likelihood, Bayes as shrinkage}
``Notice how Bayes methods are sort of like penalized likelihood methods; `shrinkage'
estimators.'' Both readings are visible in objects already on the table. \emph{Penalty}: by
\xref{2.9.1}, maximizing the posterior means maximizing
\[ \log \pi(\theta\mid x) \;=\; \ell(\theta) + \log\pi(\theta) + \text{const} : \]
the log-likelihood \emph{plus a penalty term} that docks points for straying where the prior
is thin. A normal prior contributes $\log\pi(\theta)=-(\theta-m)^2/(2\tau^2)+\text{const}$,
a quadratic penalty on distance from $m$---the ridge-regression form, with $1/\tau^2$ as
the tuning knob. \emph{Shrinkage}: the posterior-mean formulas of \S\,\xref{2.9} do not
merely resemble shrinking---they \emph{are} affine maps pulling the data's guess partway
toward the prior's, by factors like $\tau^2/(\sigma^2+\tau^2)$ and $n/(\alpha+\beta+n)$.
The deck's hedge (``sort of like'') is apt: the posterior \emph{mean} and the posterior
\emph{maximizer} are different summaries of the same posterior---they coincide in the
normal--normal case and differ in general---but the family resemblance among Bayes,
penalization, and shrinkage is structural, not cosmetic.
\end{ObjL}

\begin{ObjL}{Remark}{2.10.2}{The deck's big-data point}
The slide closes: ``At the heart of first-order big data: you cannot use the data to know
what to shrink towards; but you can use the data to know how much to shrink towards it.''
Unpacked: the shrinkage \emph{target} ($m$; the prior mean) plays the role of outside
knowledge---estimating it from the same data that will then be shrunk toward it is
circular, at least to first order. But the shrinkage \emph{amount}---the weight
$\tau^2/(\sigma^2+\tau^2)$---is a different kind of quantity: with many parallel problems
(many arms, many genes, many small areas), the spread of the observed estimates around the
target reveals $\tau^2$, and the data can tune their own shrinkage. Estimating the prior's
parameters from the ensemble is called \emph{empirical Bayes}, and the discovery that such
procedures can beat the obvious unbiased one at its own game (Stein's phenomenon) is one of
the subject's genuine surprises---here, only a signpost at the edge of the course.
\end{ObjL}

% ============================================================
\sub{2.11}{Looking Ahead: Sufficiency}

\begin{ObjL}{Concept}{2.11.1}{``If we are done by Thursday''}
The deck's final slide is a conditional promise, and a map of the next idea: \emph{sufficient
statistics}---summaries of the data that lose nothing about the parameter (note how every
example this week funneled the whole sample through $\bar X$ or a total; Fisher's ``reduction
of data,'' the epigraph, is this program); \emph{conditioning on sufficient statistics},
whose payoffs the slide indexes only as ``Unbiasedness'' and ``Relationship to
variance''---the mathematics to come: conditioning preserves unbiasedness and never
increases variance, so any unbiased estimator can be upgraded by averaging it over the
sufficient summary; and \emph{complete} sufficient statistics, whose ``implications for uniqueness''
deliver the unique uniformly optimal unbiased estimator that \xref{2.1.7} advertised. The
deck defers the mathematics, and so do these notes: the objects above are next week's, named
here so the week's arc---floor, constructor, precision, prior---is visible as one story.
\end{ObjL}

% ============================================================
\lectureappendix

The deck names its computations and leaves them to the board and to us; the longer ones are
worked here, every step shown, each keyed to the object it supports. \appref{2.A.1} settles
the sample variance's bias (\xref{2.1.6}, \xref{2.6.3}); \appref{2.A.2} prices nuisance
parameters (\xref{2.2.5}); \appref{2.A.3} computes the exponential rate's bias and the
variance of its repair (\xref{2.5.4}); \appref{2.A.4} completes the square behind the
normal--normal posterior (\xref{2.9.4}); and \appref{2.A.5} runs Fieller's quadratic through
its cases (\xref{2.7.7}).

\begin{ObjL}{Supplement}{2.A.1}{The sample variance's bias, and Bessel's correction}
For $X_1,\dots,X_n$ iid with mean $\mu$ and variance $\sigma^2$,
\[ \E\Big\{ \frac1n \sum_{i=1}^n (X_i-\bar X)^2 \Big\} \;=\; \frac{n-1}{n}\,\sigma^2 ,
   \qquad\text{so}\qquad
   \E\Big\{ \frac{1}{n-1} \sum_{i=1}^n (X_i-\bar X)^2 \Big\} \;=\; \sigma^2 . \]
\end{ObjL}
\begin{Prf}{2.A.1}{Proof of Bessel's Correction by Adding and Subtracting the True Mean.}{}
\item Route each deviation through $\mu$: $X_i-\bar X=(X_i-\mu)-(\bar X-\mu)$, square, and
      sum:
      \[ \sum_i (X_i-\bar X)^2 \;=\; \sum_i (X_i-\mu)^2
         \;-\; 2(\bar X-\mu)\sum_i (X_i-\mu) \;+\; n(\bar X-\mu)^2 . \]
\item The middle sum telescopes into the average: $\sum_i (X_i-\mu)=n(\bar X-\mu)$, so the
      middle term is $-2n(\bar X-\mu)^2$ and
      \[ \sum_i (X_i-\bar X)^2 \;=\; \sum_i (X_i-\mu)^2 \;-\; n(\bar X-\mu)^2 . \]
      (Measuring from the sample's own center \emph{always} removes exactly
      $n(\bar X-\mu)^2$ worth of scatter---the algebraic content of ``$\bar X$ has chased
      the data,'' \xref{2.1.6}(c).)
\item Take expectations: $\E\{(X_i-\mu)^2\}=\sigma^2$ for each $i$, and
      $\E\{(\bar X-\mu)^2\}=\Var(\bar X)=\sigma^2/n$, so
      \[ \E\Big\{\sum_i (X_i-\bar X)^2\Big\} \;=\; n\sigma^2 - n\cdot\frac{\sigma^2}{n}
         \;=\; (n-1)\,\sigma^2 . \]
\item Divide by $n$ for the plug-in's expectation, $(n-1)\sigma^2/n$; divide by $n-1$
      instead and the expectation is exactly $\sigma^2$.
\end{Prf}

\begin{ObjL}{Supplement}{2.A.2}{Nuisance parameters cost information: the $2\times2$ case}
Let $I=\begin{pmatrix} I_{11} & I_{12}\\ I_{12} & I_{22}\end{pmatrix}$ be an information
matrix (symmetric, positive definite). Then
\[ \big(I^{-1}\big)_{11} \;=\; \frac{I_{22}}{I_{11}I_{22}-I_{12}^2}
   \;\ge\; \frac{1}{I_{11}} ,
   \qquad\text{with equality if and only if } I_{12}=0 . \]
\end{ObjL}
\begin{Prf}{2.A.2}{Proof of the Nuisance Cost by the Explicit Inverse.}{}
\item The $2\times2$ inverse: with $\det I=I_{11}I_{22}-I_{12}^2>0$ (positive
      definiteness),
      \[ I^{-1} \;=\; \frac{1}{I_{11}I_{22}-I_{12}^2}
         \begin{pmatrix} I_{22} & -I_{12}\\ -I_{12} & I_{11} \end{pmatrix} ,
         \qquad\text{so}\qquad
         \big(I^{-1}\big)_{11} \;=\; \frac{I_{22}}{I_{11}I_{22}-I_{12}^2} . \]
\item Compare with $1/I_{11}$ by subtracting, over the common positive denominator:
      \[ \big(I^{-1}\big)_{11} - \frac{1}{I_{11}}
         \;=\; \frac{I_{11}I_{22} - \big(I_{11}I_{22}-I_{12}^2\big)}
         {I_{11}\big(I_{11}I_{22}-I_{12}^2\big)}
         \;=\; \frac{I_{12}^2}{I_{11}\big(I_{11}I_{22}-I_{12}^2\big)} \;\ge\; 0 , \]
      every factor in the denominator positive.
\item The difference vanishes exactly when $I_{12}=0$: the orthogonal case. Correlated
      scores mean the nuisance coordinate's uncertainty leaks into yours; the leak is the
      ratio displayed.
\end{Prf}

\begin{ObjL}{Supplement}{2.A.3}{The exponential rate: bias of $1/\bar X$, and the repaired estimator's variance}
Let $X_1,\dots,X_n$ be iid Exponential with rate $\lambda$ and let $S=\sum_i X_i$. Then for
$n\ge2$,
\[ \E^{\lambda}\Big\{\frac{n}{S}\Big\} \;=\; \frac{n\lambda}{n-1} , \]
so the MLE $\hat\lambda=1/\bar X=n/S$ is biased upward; and the unbiased repair
$(n-1)/S$ has, for $n\ge3$,
\[ \Var^{\lambda}\Big(\frac{n-1}{S}\Big) \;=\; \frac{\lambda^2}{n-2}
   \;>\; \frac{\lambda^2}{n} \;=\; \frac{1}{I_n(\lambda)} . \]
\end{ObjL}
\begin{Prf}{2.A.3}{Proof of the Bias and the Repaired Variance via the Gamma Integral.}{Uses \xref{2.5.4}.}
\item The total of $n$ iid Exponential$(\lambda)$ variables has the Gamma$(n,\lambda)$
      density $f_S(s)=\lambda^n s^{\,n-1} e^{-\lambda s}/\Gamma(n)$, $s>0$ (a standard
      fact; it can be built by induction on $n$).
\item Negative first moment, by matching the integrand back to a Gamma density of lower
      order:
      \[ \E\Big\{\frac1S\Big\}
         \;=\; \int_0^{\infty} \frac1s\cdot\frac{\lambda^n s^{\,n-1}e^{-\lambda s}}
         {\Gamma(n)}\,ds
         \;=\; \frac{\lambda^n}{\Gamma(n)} \int_0^{\infty} s^{\,n-2}e^{-\lambda s}\,ds , \]
      and the surviving integral is the Gamma integral one order down,
      $\Gamma(n-1)/\lambda^{\,n-1}$, so
      \[ \E\Big\{\frac1S\Big\}
         \;=\; \frac{\lambda^n}{\Gamma(n)}\cdot\frac{\Gamma(n-1)}{\lambda^{\,n-1}}
         \;=\; \frac{\lambda}{n-1} , \]
      using $\Gamma(n)=(n-1)\Gamma(n-1)$ and requiring $n\ge2$ for the integral to
      converge. Multiply by $n$: $\E\{n/S\}=n\lambda/(n-1)$.
\item Negative second moment, the same maneuver one order down:
      \[ \E\Big\{\frac{1}{S^2}\Big\}
         \;=\; \frac{\lambda^n}{\Gamma(n)}\cdot\frac{\Gamma(n-2)}{\lambda^{\,n-2}}
         \;=\; \frac{\lambda^2}{(n-1)(n-2)} \qquad (n\ge3) . \]
\item Hence for the estimator $(n-1)/S$ (unbiased by step (2), since
      $\E\{(n-1)/S\}=(n-1)\cdot\lambda/(n-1)=\lambda$):
      \[ \Var\Big(\frac{n-1}{S}\Big)
         \;=\; (n-1)^2\left[ \frac{\lambda^2}{(n-1)(n-2)}
         - \frac{\lambda^2}{(n-1)^2} \right]
         \;=\; \lambda^2\left[ \frac{n-1}{n-2} - 1 \right]
         \;=\; \frac{\lambda^2}{n-2} , \]
      strictly above the floor $\lambda^2/n$: in this parameterization no estimator in
      sight attains the bound---``a not sharp bound'' indeed (\xrefL{1.9.14}).
\end{Prf}

\begin{ObjL}{Supplement}{2.A.4}{The normal--normal posterior, every step}
With prior $\theta\sim N(m,\tau^2)$ and one observation $X\mid\theta\sim N(\theta,\sigma^2)$,
the posterior is $N(m^*,v)$ with
\[ v=\Big(\tfrac{1}{\tau^2}+\tfrac{1}{\sigma^2}\Big)^{-1} , \qquad
   m^{*}=v\Big(\tfrac{m}{\tau^2}+\tfrac{x}{\sigma^2}\Big)
   =\tfrac{\sigma^2 m+\tau^2 x}{\sigma^2+\tau^2} . \]
\end{ObjL}
\begin{Prf}{2.A.4}{Proof of the Posterior by Completing the Square in $\theta$.}{Uses \xref{2.9.1}.}
\item Multiply prior and likelihood and keep only $\theta$-dependent factors:
      \[ \pi(\theta\mid x) \;\propto\;
         \exp\!\left( -\frac{(\theta-m)^2}{2\tau^2} \right)
         \exp\!\left( -\frac{(x-\theta)^2}{2\sigma^2} \right) , \]
      so, combining the exponents,
      \[ \pi(\theta\mid x) \;\propto\;
         \exp\!\left( -\frac12\left[ \frac{(\theta-m)^2}{\tau^2}
         + \frac{(\theta-x)^2}{\sigma^2} \right] \right) . \]
\item Expand the bracket and collect powers of $\theta$:
      \[ \frac{(\theta-m)^2}{\tau^2}+\frac{(\theta-x)^2}{\sigma^2}
         \;=\; \theta^2\Big(\frac{1}{\tau^2}+\frac{1}{\sigma^2}\Big)
         - 2\theta\Big(\frac{m}{\tau^2}+\frac{x}{\sigma^2}\Big)
         + \Big(\frac{m^2}{\tau^2}+\frac{x^2}{\sigma^2}\Big) . \]
\item Write $\frac1v=\frac{1}{\tau^2}+\frac{1}{\sigma^2}$ and
      $m^*=v\big(\frac{m}{\tau^2}+\frac{x}{\sigma^2}\big)$, and complete the square:
      \[ \theta^2\cdot\frac1v - 2\theta\cdot\frac{m^*}{v}
         \;=\; \frac{\theta^2 - 2\theta\, m^*}{v}
         \;=\; \frac{(\theta-m^*)^2}{v} - \frac{(m^*)^2}{v} , \]
      and the terms free of $\theta$ join the normalizer.
\item The $\theta$-dependence is exactly $\exp\{-(\theta-m^*)^2/(2v)\}$: the $N(m^*,v)$
      kernel. Since the posterior is a probability distribution (\xref{2.9.1}), it \emph{is}
      $N(m^*,v)$. Clearing $v$ from $m^*$'s definition---multiply numerator and denominator
      by $\sigma^2\tau^2$:
      \[ m^* \;=\; \frac{\dfrac{m}{\tau^2}+\dfrac{x}{\sigma^2}}
         {\dfrac{1}{\tau^2}+\dfrac{1}{\sigma^2}}
         \;=\; \frac{\sigma^2 m+\tau^2 x}{\sigma^2+\tau^2} . \]
\end{Prf}

\begin{ObjL}{Supplement}{2.A.5}{Fieller's quadratic, case by case}
The set of \xref{2.7.7} is $S=\{\rho: g(\rho)\le 0\}$ where, writing
$c=z_{\alpha/2}^2\sigma^2$,
\[ g(\rho) \;=\; \Big(\bar Y^2-\frac{c}{m}\Big)\rho^2 \;-\; 2\bar X\bar Y\,\rho
   \;+\; \Big(\bar X^2-\frac{c}{n}\Big) , \]
and its shape is decided by the leading coefficient and the roots.
\end{ObjL}
\begin{Prf}{2.A.5}{Proof of the Case Analysis by Expanding the Inversion Inequality.}{Uses \xref{2.7.7}.}
\item Square out the defining inequality
      $(\bar X-\rho\bar Y)^2 \le c\,\big(\frac1n+\frac{\rho^2}{m}\big)$ and move everything
      to the left:
      \[ \bar X^2 - 2\bar X\bar Y \rho + \bar Y^2\rho^2 - \frac{c}{n}
         - \frac{c}{m}\rho^2 \;\le\; 0 , \]
      which is $g(\rho)\le0$ with $g$ as displayed.
\item The discriminant of $g$, after expanding and cancelling the $\bar X^2\bar Y^2$
      terms:
      \[ 4\bar X^2\bar Y^2 - 4\Big(\bar Y^2-\frac{c}{m}\Big)\Big(\bar X^2-\frac{c}{n}\Big)
         \;=\; 4c\left( \frac{\bar Y^2}{n} + \frac{\bar X^2}{m} - \frac{c}{nm} \right) . \]
\item In the case of a positive leading coefficient ($\bar Y^2>c/m$: the denominator mean
      is convincingly nonzero), the parabola opens upward and the discriminant is strictly
      positive: since $\bar Y^2>c/m$, its parenthesis strictly exceeds
      $\frac{c/m}{n}+\frac{\bar X^2}{m}-\frac{c}{nm}=\frac{\bar X^2}{m}\ge0$. Two real
      roots; $S$ is the genuine interval between them.
\item In the case of a negative leading coefficient ($\bar Y^2<c/m$: the data cannot rule
      out $\mu_2=0$), the parabola opens downward. If the discriminant is positive, $g\le0$
      outside the roots: $S$ is two half-lines running to $\pm\infty$. If it is negative,
      $g<0$ everywhere: $S=\mathbb{R}$---the data place no restriction on the ratio at all.
\item Coverage is untouched by the case analysis: whatever its shape, $S$ contains the true
      $\rho$ exactly when $|Z(\rho)|\le z_{\alpha/2}$, an event of probability $1-\alpha$
      (\xref{2.7.4}). When the denominator might be zero the ratio is genuinely
      ill-determined, and an honest $1-\alpha$ report says so by refusing to be an
      interval.
\end{Prf}

\ifSubfilesClassLoaded{\clearpage\objindex}{}

\end{document}
